\documentclass[11pt,letterpaper]{article}

\usepackage[margin=1in]{geometry}
\usepackage[T1]{fontenc}
\usepackage{lmodern}
\usepackage{microtype}
\usepackage{parskip}
\usepackage{enumitem}
\usepackage{booktabs}
\usepackage{longtable}
\usepackage{array}
\usepackage{titlesec}
\usepackage[hidelinks,colorlinks=true,linkcolor=black,urlcolor=black,citecolor=black]{hyperref}
\usepackage{xcolor}
\usepackage{fancyhdr}
\usepackage{amsmath}
\usepackage{amssymb}

\titleformat{\section}{\Large\bfseries}{\thesection}{1em}{}
\titleformat{\subsection}{\large\bfseries}{\thesubsection}{1em}{}

\pagestyle{fancy}
\fancyhf{}
\fancyhead[L]{\small The Theseus Guides --- III. Research Foundations}
\fancyhead[R]{\small\thepage}
\renewcommand{\headrulewidth}{0.4pt}

\setlength{\parindent}{0pt}

\newcommand{\file}[1]{\texttt{#1}}

\title{\textbf{The Theseus Guides}\\[6pt]\large III. Research Foundations:\\The Empirical Substrate of the Firm's Methods}
\date{April 2026}
\author{}

\begin{document}
\maketitle
\thispagestyle{empty}

\begin{quote}\itshape
Every method the firm deploys rests on a testable empirical claim about the geometry of meaning or the structure of contradiction. This guide catalogues those claims and the experiments that support them.
\end{quote}

\vspace{1em}
\tableofcontents
\newpage

% =======================================================================
\section{Why this guide exists}
% =======================================================================

Guide II gave the firm's formal commitments: five criteria that any reliable method must satisfy. A pure philosophical commitment, however rigorous, is insufficient for engineering. The firm also needs empirical foundations: specific, testable claims about how meaning and contradiction behave in the mathematical spaces modern language models produce. These claims are what make the coherence engine, the contradiction detector, and the adversarial generator possible at all.

The research directory for this empirical work is \file{ideologicalOntology/} in the repository. This guide summarizes the four principal research programmes found there: the Quintin Hypothesis, the Embedding Geometry Conjecture, the Contradiction Geometry experiment suite, and the Reverse Marxism ideology-flipping framework. Each is an independently-testable claim; together they constitute the firm's wager that meaning has enough geometric structure to support disciplined truth-finding at scale.

% =======================================================================
\section{The Quintin Hypothesis}
% =======================================================================

\textbf{Statement.} Logical coherence --- the absence of contradiction --- is a \emph{geometric} property detectable in large language model embedding space, and claim-sets whose geometric coherence is higher predict real-world outcomes better than claim-sets whose coherence is lower.

\textbf{Why it matters.} The hypothesis is load-bearing for the entire firm. If coherence is \emph{not} geometrically detectable, the coherence engine is not measuring anything real. If geometrically detected coherence does \emph{not} predict outcomes, the firm's wager collapses: no amount of internal consistency would translate into external truth. The hypothesis is the bridge between the second-order methodological commitments of Guide II and the first-order engineering of Guides IV--V.

\textbf{Status.} The hypothesis has two components, empirically separable:

\begin{enumerate}[leftmargin=1.5em,itemsep=0.2em]
\item \emph{Detectability.} Contradictions manifest as a specific geometric configuration in embedding space. This has been validated (see the Embedding Geometry Conjecture, below), and is the substrate on which the six-layer coherence engine is built.
\item \emph{Predictivity.} Higher geometric coherence correlates with better out-of-sample predictions. This is what the Calibration Scoreboard tests continuously; it is a \emph{progressive} (in the Lakatosian sense) hypothesis, making specific predictions that are being scored over time.
\end{enumerate}

The firm does not assert the hypothesis as settled. It asserts the hypothesis as \emph{testable}, and it operates the apparatus that tests it. This distinction is the meta-method applied to the firm's own most ambitious claim.

% =======================================================================
\section{The Embedding Geometry Conjecture}
% =======================================================================

\textbf{Statement (from \file{ideologicalOntology/Embedding\_Geometry\_Conjecture/README.md}).}

\begin{quote}
LLMs represent every concept as a point in trillion-dimensional space. Logical and semantic relationships become geometric relationships. \textbf{Conjecture:} A contradiction between two statements manifests as a specific geometric configuration --- vectors pointing in incompatible directions.
\end{quote}

The conjecture refines the na\"ive ``contradiction $=$ opposite vectors'' picture, which is false. Contradictions in practice have \emph{high} cosine similarity because the two statements share topic. ``The sky is blue'' and ``The sky is not blue'' live near each other in embedding space, as do ``Capitalism creates prosperity'' and ``Capitalism creates poverty.'' Raw cosine similarity is useless for contradiction detection. This is the \emph{Cosine Paradox} identified in the research directory's \file{Coherence\_Formula\_Walkthrough.pdf}.

\subsection{The seven tests}

The validation experiment suite at \file{ideologicalOntology/Embedding\_Geometry\_Conjecture/} runs seven tests, whose design and expected outcomes are documented in the README:

\begin{enumerate}[leftmargin=1.5em,itemsep=0.15em]
\item \textbf{Cosine Similarity Distributions.} Do contradictions point in opposite directions? Spoiler: no.
\item \textbf{Difference Vector Sparsity.} Is the contradiction signal concentrated in a small number of dimensions?
\item \textbf{Self-Consistency.} Is there a universal ``contradiction direction'' vector $\hat{c}$?
\item \textbf{Linear Separability.} Can a linear classifier distinguish contradiction from entailment from neutral?
\item \textbf{PCA of Difference Vectors.} What is the intrinsic dimensionality of the contradiction subspace?
\item \textbf{Signal Concentration.} Which specific dimensions carry the contradiction signal?
\item \textbf{Hard Cases.} Do minimal pairs (``X is Y'' vs.\ ``X is not Y'') break under the tested method?
\end{enumerate}

\subsection{Key findings}

The synthetic and real-model experiments (files \file{simulation\_synthetic.py}, \file{experiment\_real\_models.py}, and \file{experiment\_contradiction\_direction.py}) establish, jointly with the literature the experiments draw on (Marks \& Tegmark 2023 on truth as a linear direction; the SparseCL work on contradiction subspaces; the ``Semantics at an Angle'' 2025 work on directional semantics), the following:

\textbf{Contradictions have high cosine similarity.} This is a published empirical fact in the literature and is reproduced in the synthetic simulation; it is the reason the na\"ive approach fails.

\textbf{The contradiction signal is sparse.} The difference vector $v_1 - v_2$ between a contradictory pair is concentrated in a small fraction (the simulation models this as $\sim$14\%) of the embedding dimensions. Hoyer sparsity of the difference vector is therefore a useful discriminator.

\textbf{A universal contradiction direction exists.} A linear probe trained on general-domain contradictory pairs generalizes to philosophy-domain and business-domain pairs, suggesting that the direction $\hat{c}$ separating contradictory from non-contradictory difference vectors is not strongly domain-specific.

\textbf{Linear classifiers outperform cosine-only baselines.} A logistic regression or linear SVM in the difference-vector space reliably distinguishes the three classes (contradiction, entailment, neutral), with performance substantially above the cosine-only baseline. Combining the cosine feature with the sparsity feature and the linear-probe projection is the strongest detector.

\subsection{How the conjecture enters production code}

The findings are operationalized in two principal modules:

\file{noosphere/noosphere/geometry.py} implements contradiction-geometry detection via Householder reflection and difference-vector sparsity per the validated conjecture. It is consumed by Layer 4 of the coherence engine.

\file{noosphere/noosphere/coherence/geometry.py} is the coherence-engine wrapper: it provides the per-pair geometric verdict that enters the six-layer aggregator alongside the NLI layer, the argumentation layer, the probabilistic layer, the information-theoretic layer, and the LLM-as-judge layer.

In both cases the conjecture's role is the same: it provides a second, orthogonal read on contradiction. A pair that the NLI head calls contradictory but that the geometry layer calls neutral triggers the aggregator's \texttt{unresolved} output, not a point verdict. The design is that no single layer wins; cross-layer agreement is the unit of inference, and disagreement is a first-class signal.

% =======================================================================
\section{The Contradiction Geometry experiment suite}
% =======================================================================

The \file{ideologicalOntology/Contradiction\_Geometry/} directory holds the broader experimental programme that the Embedding Geometry Conjecture is part of. It contains nine completed experiments documented in \file{Methods\_and\_Experiments\_Reference.md}, covering:

\begin{itemize}[leftmargin=1.5em,itemsep=0.1em]
\item Contradiction detection across model families (MiniLM, MPNet, BGE-large, OpenAI embeddings).
\item Ontological mapping: building directed graphs whose nodes are atomic claims and whose edges are typed relations (supports / contradicts / refines / instantiates).
\item Formal mathematical framework for the coherence score itself, documented in \file{Formal\_Mathematical\_Framework.pdf}.
\item Three-dimensional visualizations of contradiction reflections (\file{contradiction\_reflection\_3d.html}).
\end{itemize}

The experiments serve a specific epistemic role: they are the \emph{severity apparatus} (in Mayo's sense) for the Embedding Geometry Conjecture. They would detect the conjecture's falsity under plausible failure modes: model-specific idiosyncrasies, cross-domain transfer failures, negation-style robustness failures, minimal-pair failures. That the conjecture has survived this apparatus is what justifies its use in production.

The detailed coherence-formula walkthrough at \file{Coherence\_Formula\_Walkthrough.pdf} documents how the numeric coherence score in production is computed from the experiment-validated primitives. It is the bridge between the research directory and the coherence engine in \file{noosphere/noosphere/coherence/}.

% =======================================================================
\section{Reverse Marxism: ideology flipping by concept-axis reflection}
% =======================================================================

The Reverse Marxism research (\file{ideologicalOntology/Reverse\_Marxism\_Ideology\_Flipping\_Research.md}) is a distinct but complementary line of empirical work. It asks: \emph{what happens when one reflects an entire ideological structure across the hyperplane defined by the concept that structure is organized around?}

The technical core is simple. Given an embedding vector $v$ and a unit vector $\hat{a}$ that represents a concept axis (``class,'' ``liberty,'' ``efficiency,'' etc.), the Householder reflection across the hyperplane orthogonal to $\hat{a}$ is
\[
v' = v - 2\,(v \cdot \hat{a})\,\hat{a}.
\]
This reflection is an involution: applying it twice returns $v$. It preserves all components of $v$ orthogonal to $\hat{a}$ and inverts exactly the $\hat{a}$-component. Sentences in an ideological text that are entirely about the concept axis are fully inverted; sentences that are orthogonal to the concept axis are preserved; sentences that lie at an angle are partially transformed in proportion to their projection onto $\hat{a}$.

\subsection{Why reflection rather than negation}

Simple negation produces nonsense: ``capitalism exploits workers'' becomes ``capitalism does not exploit workers,'' which is a denial rather than a constructive counter-framework. Reflection preserves the \emph{structure} of the argument while inverting its \emph{orientation}. Marx's analysis of how value flows through the economy --- the \emph{fact} that there are systematic flows between classes --- is preserved under reflection. What changes is the \emph{direction} of the flow: instead of ``value extracted from worker to capitalist (exploitation),'' the reflected structure describes ``value delivered from capitalist to worker (enrichment).'' The reflected text reads as a systematic pro-capitalist argument with the structural rigor of Marx, because it inherits that structural rigor under the reflection operation.

\subsection{The hypothesis being tested}

The research hypothesis is stated in the markdown file explicitly:

\begin{quote}
The reflected structure IS capitalism's own self-justification --- the argument that capitalism makes for itself, articulated with the same depth and systematicity that Marx brought to his critique.
\end{quote}

This is a surprising, testable, and falsifiable empirical claim. It would be falsified if the reflected text reads as incoherent noise, or if the nearest corpus sentences to the reflected vectors form an ideologically random collection, or if domain experts asked to rate the reflected text for coherence and ideological recognizability gave it low scores. Experiments 1 and 2 in the research programme are designed to test exactly this.

\subsection{Why the firm cares about Reverse Marxism}

Reverse Marxism is a \emph{method for generating the strongest counter-argument to any position}. That role is load-bearing in two places in the firm's engineering:

\textbf{Adversarial challenge generation.} The adversarial subsystem (\file{noosphere/noosphere/adversarial.py}) generates structured objections to firm-tier conclusions. The Reverse Marxism reflection is the formal basis for ``mechanically generate the strongest counter-argument,'' which is the severity requirement from Guide II made executable. Without a principled way to generate the strongest counter-position, adversarial challenge degenerates into whatever the LLM happens to produce, which is at best a weak steelman.

\textbf{Dialectical Crucible (Guide VI, Project 2).} The planned computational-dialectics module takes this further: every investment thesis, every published conclusion, every prediction is run through thesis $\to$ antithesis-by-reflection $\to$ synthesis-by-methodological-quality. The output is not ``which side is right'' but ``which side reasons better.'' This is Hegelian dialectics reconstructed as computation, with the reflection operation standing in for the traditional Hegelian negation.

\subsection{The broader research programme}

The Reverse Marxism file sketches a research programme with increasing scale:

\begin{itemize}[leftmargin=1.5em,itemsep=0.15em]
\item \textbf{Experiment 1} (proof of concept): reflect 50 key sentences from the \textit{Communist Manifesto} across the ``class'' axis; evaluate the resulting text for coherence and ideological recognizability.
\item \textbf{Experiment 2} (full scale): reflect the entirety of \textit{Capital, Vol.~1} (\textasciitilde30,000 sentences) across a validated class axis, with a 2M-sentence reference corpus (Wikipedia plus Smith, Friedman, Rand, Mises, Hayek, Rothbard, Sowell, Thiel, Taleb).
\item \textbf{Generalization experiments}: reflect other ideological systems across their organizing axes --- libertarianism across ``coercion,'' progressivism across ``equality,'' conservatism across ``tradition.'' Each reflection is a hypothesis test: does the reflected structure correspond to a recognizable real-world ideology?
\end{itemize}

The generalization experiments are where the programme becomes most ambitious. If reflections across the right axes reliably produce recognizable real-world ideologies, the implication is that the space of coherent ideological systems is a \emph{small, geometrically-structured set}, not a vast combinatorial landscape --- and that ideologies are related by reflections across a limited number of organizing axes. This is a claim about the structure of political thought that is unusual, testable, and would matter if true.

% =======================================================================
\section{Domain-relative coherence}
% =======================================================================

A companion research document, \file{ideologicalOntology/Domain\_Relative\_Coherence.md}, extends the framework in a direction that matters for how the coherence engine is used in practice. Its argument is that coherence should be measured \emph{within a domain} rather than against a global set of societal premises, because societal premises are themselves incoherent (the document reports a measured overall coherence of 0.321 on a 47-premise database of operative liberal-democratic beliefs, with 12 cross-domain tensions). Measuring against an incoherent yardstick gives measurements of doubtful semantic content.

The document identifies five candidate definitions of ``domain'' --- topical, epistemic, normative, ontological, functional --- and argues that a domain is best understood as an \emph{emergent property of the argument's own structure}, defined by the premises it invokes, the inference patterns it uses, the entities it presupposes, and the standards of evidence it accepts.

This matters architecturally because it explains why the firm's coherence engine does not pursue a single global coherence score, why the ontology module (\file{noosphere/noosphere/ontology.py}) builds growing topic ontologies rather than imposing a fixed taxonomy, and why the meta-analysis gate (\file{noosphere/noosphere/meta\_analysis.py}) applies its five criteria \emph{relative to} the candidate conclusion's domain of operation rather than against a universal benchmark.

% =======================================================================
\section{The written catalogue}
% =======================================================================

The firm's research documentation lives in two complementary places. The empirical experiments are in \file{ideologicalOntology/}, summarized above. The written work that motivates and constrains the software is in \file{docs/} as PDFs; the most relevant for understanding the system as a whole are:

\begin{longtable}{p{5cm} p{9.5cm}}
\toprule
\textbf{Document} & \textbf{What it covers} \\
\midrule
\endhead
\file{Methodological\_Reorientation.pdf} & The pivot. Defines why Theseus's product is methodology, not claims. Governs every downstream architectural choice. \\
\file{The\_Meta\_Method.pdf}             & The firm's methodological commitments in detail. \\
\file{Product\_Description.pdf}          & How the three user-facing softwares fit together. \\
\file{Noosphere\_Project\_Status.pdf}    & Engineering status snapshot. \\
\file{Server\_Architecture.pdf}          & Deployment architecture and the public/private split. \\
\file{Geometry\_of\_Unresolution.pdf}    & Mathematical treatment of epistemic unresolvability; the formal substrate behind \texttt{unresolved} as a first-class output. \\
\file{Podcast\_Infrastructure\_Research.pdf} & Research motivating the Dialectic audio pipeline and Noosphere's ingester. \\
\file{Transcript\_Extraction\_Research.pdf}  & Research motivating the chunk-to-claim extraction pipeline. \\
\file{Founders\_Reading\_and\_Research.pdf}  & Reading list and background material. \\
\file{Founders\_Questions.pdf}           & Active open questions the firm is tracking. \\
\file{Performance\_Synthesis.pdf}        & Method performance and synthesis strategies. \\
\file{Voices\_Method.pdf}               & Voice-decomposition methodology for separating individual founder positions. \\
\bottomrule
\end{longtable}

Every one of these is either cited by an implementation or constrains an implementation decision. They are not vestigial documentation; they are the firm's intellectual justification for why the software is shaped the way it is, and in the event of ambiguity between what the code does and what the documents specify, the documents have been the source of truth for round-3 architectural decisions.

\end{document}
